\subsection*{Writing a fact beside an attested historical form}
Asked for the original meaning and bronze form of the pictured headword for ear, a 28,032-parameter writer produced a complete Chinese sentence containing the English gloss \emph{ear} and the corresponding Hanziyuan inscription (Fig.~\ref{fig:pictured-entries}). A second question about the same entry elicited its pronunciation and a different historical form. Both answers matched the prepared reference at every pixel. Across 142 cards from 284 held-out character families, 1,134 of 1,136 complete answers were exact. All eight requests were correct on 141 cards. Reassigning questions to different answers within each training card gave zero correct development sentences. An authored parser that knows the two question frames and the answer convention solved all 1,136.

When corresponding values changed places between the two entries and each question stayed fixed, 1,134 answers followed the revised evidence. A new question frame retaining the printed field names gave 1,132 exact answers. Replacing the field names with synonyms reduced the score to 225, even though 1,129 responses still ended with the requested form and punctuation. In a further diagnostic, equal-length synonyms for meaning and pronunciation gave exactly the same input features for all 568 pairs requiring different answers. The recurrent decoder starts in the same state on both sides of each pair. Identical candidate images and features give identical greedy decisions at each subsequent step, so even optimal weights under this representation could resolve zero pairs in both directions. The model learned to assemble the requested evidence across families. Extending this result to unfamiliar expressions requires a reading pathway that preserves the difference between the questions.
